
\chapter*{Frontispiece}
\addcontentsline{toc}{chapter}{Frontispiece}
\label{ch:frontispiece}

The Igusa cusp form \(\Delta_5\) has five readings.  The genus-two
automorphic section, the denominator of a generalized
Borcherds--Kac--Moody superalgebra on \(\La^{2,1}_{II}\), and the
singular theta lift of the weak Jacobi form \(\phi_{0,1}\) are
theorem-level readings.  The protected Pfaffian is a retained-data
reading of the chosen \(E_3\)-prefactorization datum
\[
  \mathfrak A_{K3\times E}
  \;\in\; E_3\text{-}\mathrm{HolFA}(K3\times E),
\]
and the mirror discriminant is a conjectural period-side reading on the
rank-\(3\) Mukai-invariant \(K3\times E\) component.

\paragraph{The five views.}
\begin{enumerate}
\item[\textcircled{1}]
\textsl{Automorphic.}
The Borcherds--Gritsenko--Nikulin section
\[
  \mathcal D_X \;=\; \Delta_5
  \;\in\; H^0\bigl(\Sp(2,\Z)\backslash\mathcal H_2,\;\mathcal L^{\otimes 5}
  \otimes\nu_{\Delta_5}\bigr)
\]
of the weight-\(5\) automorphic line on Siegel upper-half space, with
multiplier \(\nu_{\Delta_5}\) of \cite{MAASS:1, Borcherds95}.

\item[\textcircled{2}]
\textsl{Borcherds--Kac--Moody denominator.}
The denominator of the generalized BKM Lie superalgebra
\(\mathfrak g_{\Delta_5}\) on the rank-three Lorentzian root lattice
\(\La^{2,1}_{II}\subset\La^{2,1}=\HypPlan\oplus\langle 2\rangle\) of
signature \((2,1)\), the index-four sublattice with even \(f_{\pm 2}\)
coordinates, abstractly \(U(4)\oplus\langle 2\rangle\) (discriminant
\(-32\), the index-four scaling of the hyperbolic block).  The ambient
lattice \(\La^{2,1}\) has basis \((f_2,f_3,f_{-2})\) with
\((f_2,f_{-2})=-1\) and \((f_3,f_3)=2\); the sublattice
\(\La^{2,1}_{II}\) is generated by \(2f_2,f_3,2f_{-2}\).  The type-II
Cartan matrix on simple roots
\(\delta_1 = 2f_2 - f_3\), \(\delta_2 = 2f_{-2} - f_3\),
\(\delta_3 = f_3\) is, with \(\mathbf{J}_3\) the \(3\times 3\)
all-ones matrix,
\((\delta_i, \delta_j) = 4\Id_3 - 2\mathbf{J}_3 = \bigl(\begin{smallmatrix}
2 & -2 & -2\\ -2 & 2 & -2\\ -2 & -2 & 2 \end{smallmatrix}\bigr)\),
\[
  \operatorname{den}(\mathfrak g_{\Delta_5})
  \;=\; 64^{-1}\,\Delta_5(2Z),
\]
with signed root supermultiplicities
\(\smult\,\mathfrak g_{\Delta_5,\alpha(n,l,m)} = f(nm,l)\) on the
active support \(\Gamma_{\mathrm{act}} = \{(n,l,m)\in\Gamma_{\mathrm{eff}}
\mid f(nm,l)\ne 0\}\), with
\(\alpha(n,l,m) = 2nf_2 - lf_3 + 2mf_{-2}\)
\cite{Borcherds95, GN, GNPartI, GNII}.

\item[\textcircled{3}]
\textsl{Compact Pfaffian.}
The retained Pfaffian line of the \(K3\times E\) chiral
\(E_3\)-prefactorization input, presented at finite stage \(R\) by the
Dirac--Igusa pro-object
\[
  \mathfrak D^{DI}_X \;=\;
  \varprojlim_R\;
  \bigl(\mathcal A^{E_3}_{X,R},\,F^{\mathrm{hyb}}_{X,R},\,\Gamma_{X,R},
  \,\Pi_X,\,\Phi_R,\,o_R,\,H_R,\,P^{\Pi}_R,\,C_{X,R},\,
  \Theta_{\mathrm{Kos},R},\,\mathcal L_{\mathrm{Pf},R},\,\mathrm{pf}_{X,R},
  \,\epsilon_{o,R},\,\mathrm{Rec}_R\bigr)
\]
(the chosen \(E_3\)-prefactorization datum \(\mathcal A^{E_3}\), the hybrid factorization
carrier \(F^{\mathrm{hyb}}\) of \(\mathrm{Ran}^{\mathrm{hyb}}(E)\), the
Mukai charge lattice \(\Gamma\), the Gram projection \(\Pi_X\), the
vanishing-cycle coefficient sheaf \(\Phi\), the orientation \(o\), the Hall datum
\(H\), the Gram-descended primitive \(P^\Pi\), the source coalgebra
\(C\), the Koszul cobar \(\Theta_{\mathrm{Kos}}\), the Pfaffian line
\(\mathcal L_{\mathrm{Pf}}\), the Pfaffian section \(\mathrm{pf}\), the
orientation character \(\epsilon_o\), and the recognition map
\(\mathrm{Rec}\); Chapter~5 specifies the finite-stage data and their
residuals, with the \(E_3\), wrapped, Koszul, and recognition lanes
retained or conditional exactly as stated there).
The obstruction \((D0)\) is reduced in
Appendix~\ref{app:obstruction-discharges} to the retained D0-HN datum;
\((O1)\), \((O1)^+\), and \((O2)\) are reduced there to retained quotient-orientation, Weyl
transport, and type-II wall-atlas data.
Together these data supply the Pfaffian and orientation part of the
Pfaffian--Dirac theorem
\[
  \mathrm{Pf}_{\mathrm{prot}}(\mathfrak D^{DI}_X) \;=\; \Delta_5,
  \qquad
  \epsilon_o \;=\; \nu_{\Delta_5}.
\]
Here \(\epsilon_o\) is the orientation character of the Pfaffian line on
\(W^{(2)}(\La^{2,1}_{II})\), and \(\nu_{\Delta_5}\) is the Maass
multiplier of \cite{MAASS:1} on \(\Sp_4(\Z)\), pulled back along the
exterior-square isogeny \(\PSp_4(\R)\simeq\SO(3,2)_+\) and the
inclusion \(\La^{2,1}_{II}\hookrightarrow\La^{3,2}\) to a character on
\(W^{(2)}(\La^{2,1}_{II})\) taking value \(-1\) on each of the three
simple type-II reflections \(s_{\delta_i}\); see
Theorem~\ref{thm:pentadic-pfaffian-dirac}.  The symbol
\(\mathcal P_X\) denotes the protected primitive Lie superalgebra
\(P_X^{\Pi,\mathrm{red}}\) of Chapter~5.6 (the \(H^0\) of the protected
primitives of the hybrid factorization sheaf
\(\mathcal F_X\) after Gram-projection pushforward
\(\overline\Pi^\Theta_{X,*}\) and Hopf-radical quotient
\(\overline{\operatorname{Rad}\langle\,,\,\rangle_X}\); the three
notations \(\mathcal P_X = P_X^{\Pi,\mathrm{red}} = P_X\) denote one
object, abbreviated \(P_X\) where the construction data are not in
play).  The Lie-superalgebra
recognition
\[
  \mathcal P_X \;\cong\; \mathfrak g_{\Delta_5}
\]
is the separate primitive-recognition theorem of Chapter~\ref{ch:pfaffian-dirac},
conditional on the finite compact-source datum \((S1)\)--\((S10)\).

\item[\textcircled{4}]
\textsl{Mirror discriminant.}
\textsl{Conjectural.} On the rank-\(3\) Mukai-invariant
\((K3\times E)\)-period component, the discriminant section is
conjecturally \(\Delta_5^{-2}\); its polar divisor is twice the
reduced mirror discriminant, and its reduced automorphic Heegner support
is the Humbert divisor where \(\Delta_5\) vanishes
(Chapter~\ref{ch:view4-mirror}).

\item[\textcircled{5}]
\textsl{Singular theta lift.}
The theta decomposition \(\phi^{\mathrm{Weil}}_{0,1}\) of the
weight-zero index-one weak Jacobi form \(\phi_{0,1}\) has weight
\(-1/2\) for the Weil representation of \(\La^{3,2}\); its Borcherds
regularized theta lift is an orthogonal automorphic form, and the
exceptional isogeny \(\PSp_4(\R)\simeq\SO(3,2)_+\) identifies
that form with the genus-two Siegel cusp form
\[
  \Theta_{\mathrm{sing}}^{\mathrm{Bor}}(\phi^{\mathrm{Weil}}_{0,1}) \;=\; \Delta_5.
\]
This is the Borcherds--Gritsenko--Nikulin theta identity of
\cite{Borcherds95, B, GN}.
\end{enumerate}

\paragraph{The cross-level identities.}
The three edges
\(\textcircled{1}\!\leftrightarrow\!\textcircled{2}\),
\(\textcircled{2}\!\leftrightarrow\!\textcircled{5}\),
\(\textcircled{1}\!\leftrightarrow\!\textcircled{5}\) are theorems
(Borcherds 1995; Gritsenko--Nikulin 1998); they identify the
automorphic, BKM, and regularized-theta descriptions. The edge
\(\textcircled{2}\!\leftrightarrow\!\textcircled{3}\) is the
conditional Pfaffian--Dirac theorem after \((D0)\), \((O1)\),
\((O1)^+\), and the retained type-II wall-atlas datum
\((O2_{\mathrm{atlas}})\), together with the finite geometric Pfaffian
datum \((P_{\mathrm{fin}})\), on the \(K3\times E\) specialization in
Appendix~\ref{app:obstruction-discharges}
(Theorems~\ref{thm:G-D0}, \ref{thm:G-O1}, \ref{thm:G-O1-plus},
\ref{thm:G-O2-orbit-transport}), making
\(\textcircled{2}\!\leftrightarrow\!\textcircled{3}\) a theorem at the
Pfaffian and orientation level on \(K3\times E\); primitive
Lie-superalgebra recognition remains conditional on the cofinal
finite source-recognition datum and the finite Koszul comparison
datum of Definition~\ref{def:finite-koszul-comparison-datum}.  This
edge relates the BKM denominator algebra to the Dirac--Igusa Pfaffian.
The edge
\(\textcircled{1}\!\leftrightarrow\!\textcircled{3}\) is the
two-step composite \(\mathcal D_X = \Delta_5 = \mathrm{Pf}_{\mathrm{prot}}
(\mathfrak D^{DI}_X)\), with the orientation match
\(\epsilon_o = \nu_{\Delta_5}\) on \(W^{(2)}(\La^{2,1}_{II})\) supplying
the multiplier.  The protected primitive Lie superalgebra \(\mathcal P_X\)
recovers the BKM target precisely when the source representatives,
relations, pairings, radical quotient, PBW comparison, and full parity
dimensions of \((S1)\)--\((S10)\) are supplied.

\paragraph{The level-\(n\) scalar shadow.}
On the closed cobordism \(K3\times E\mapsto\mathrm{pt}\),
\[
  Z^{K3\times E}_{\mathrm{BPS}}
  \;=\; \bigl(\Delta_{10}\bigr)^{-1}
  \;=\; \Delta_5^{-2},
  \qquad
  \Delta_{10} \;=\; \Delta_5^2
  \;\in\; M_{10}\bigl(\Sp_4(\Z)\bigr).
\]
Here \(\Delta_{10} = \Delta_5^2\) is the theta-normalized weight-ten
cusp form on \(\Sp_4(\Z)\), equivalently \(\Phi_{10}^{\mathrm{un}} =
\chi_{10} = \Delta_5^2\) in the unscaled Igusa convention
(\cite{Igusa1962}; the multiplier \(\nu_{\Delta_5}\) squares to
trivial \cite{MAASS:1}); the Oberdieck--Pixton-monic normalization
\(\chi_{10}^{\mathrm{OP}} := D_5^2 = 4096^{-1}\Delta_5^2\) (with
\(D_5 := 64^{-1}\Delta_5\) the Borcherds-monic form so that
\([q^{1/2}r^{1/2}s^{1/2}]D_5 = 1\)) gives the equivalent form
\(Z^X_{\mathrm{OP}} = -(\chi_{10}^{\mathrm{OP}})^{-1} =
-4096\,\Delta_5^{-2}\), where \(4096 = 64^2\) is the squared
theta-leading constant
\([q^{1/2}r^{1/2}s^{1/2}]\Delta_5 = 64\) of the Igusa
form and the minus sign is the OP scalar branch convention
\cite{OB, OPand}.
\(Z^{K3\times E}_{\mathrm{BPS}}\) is the level-\(n\) protected trace of
the retained \(K3\times E\) data.  A three-dimensional gravity-line
interpretation with internal compact factor \(K3\times E\) requires the
Hall--Borcherds residual of Vol II and is not asserted here.

\paragraph{The \(\kappa\)-tuple of \(K3\times E\).}
The retained \(K3\times E\) data carry the five-coordinate
\(\kappa\)-tuple
\[
  \kappa_{K3\times E}
  \;=\;
  \bigl(\kappa_{\mathrm{cat}},\;\kappa^{\mathrm{Hodge}}_{\mathrm{ch}},
  \;\kappa^{\mathrm{Heis}}_{\mathrm{ch}},\;\kappa_{\mathrm{BKM}},
  \;\kappa_{\mathrm{fiber}}\bigr)
  \;=\; (0,\,0,\,3,\,5,\,24).
\]
The categorical Euler is \(\kappa_{\mathrm{cat}} =
\chi(\mathcal O_{K3\times E}) = \chi(\mathcal O_{K3})\,\chi(\mathcal O_E)
= 2\cdot 0 = 0\). The Hodge supertrace
\(\kappa^{\mathrm{Hodge}}_{\mathrm{ch}}(K3\times E) =
\sum_q (-1)^q h^{0,q}(K3\times E) = 0\). The Heisenberg rank
\(\kappa^{\mathrm{Heis}}_{\mathrm{ch}} = 3\) is the rank of the
type-II Heisenberg--Cartan obtained after the normal-ordered Gram
projection \(\Pi_X\) to \(\La^{2,1}_{II}\).  It is not the rank of the
algebraic Mukai lattice \(H^0(S,\Z)\oplus\NS(S)\oplus H^4(S,\Z)\);
for the \(U\oplus\langle-2\rangle\) polarisation that lattice has rank
five, while the \(\kappa^{\mathrm{Heis}}_{\mathrm{ch}}\)-coordinate
counts only the three Gram coordinates which enter the type-II BKM
Cartan. The BKM coordinate
\[
  \kappa_{\mathrm{BKM}}(\Phi_N) \;=\; c_N(0)/2
\]
(Borcherds 1995; Gritsenko 1999) reads, at the paramodular input
\(N=1\), \(c_1(0) = 10\), giving
\(\kappa_{\mathrm{BKM}}(\Delta_5) = 5\). The fibre coordinate
\(\kappa_{\mathrm{fiber}} = \chi_{\mathrm{top}}(K3) = 24\). The additive
identity
\(\kappa_{\mathrm{BKM}} = \kappa^{\mathrm{Hodge}}_{\mathrm{ch}}
+ \chi(\mathcal O_{\mathrm{fiber}})\) fails at every
\(N\in\{1,2,3,4,6\}\): the Lefschetz Euler \(\chi_{\mathrm{top}}(K3) =
24\) and the holomorphic Euler \(\chi(\mathcal O_{K3}) = 2\) are
distinct invariants of the K3 fibre, and \(\kappa_{\mathrm{BKM}}\) sees
the \(c_N(0)/2\) data of the Borcherds input rather than either of
them. Each subscript on \(\kappa\) records a distinct invariant;
the unsubscripted handle leaves the categorical dimension unspecified.

\chapter{Pentadic \texorpdfstring{\(\Delta_5\)}{Delta\_5}}
\label{ch:pentadic}

\section{The retained \texorpdfstring{\(E_3\)}{E\_3}-object and the five
readings of \texorpdfstring{\(\Delta_5\)}{Delta\_5}}
\label{sec:retained-e3-five-readings}

The five readings attached to \(\Delta_5\) have different logical
status.  The automorphic, BKM, and theta-lift readings are theorems;
the Pfaffian reading is relative to retained \(K3\times E\) data; the
mirror reading is conjectural.  The
trivial-canonical threefold \(K3\times E\) supplies the ambient CY-\(3\)
input for a retained holomorphic prefactorization datum in the
Kontsevich--Tamarkin and Costello--Gwilliam--Li lineage
\cite{KontsevichOperadsMotives, CostelloGwilliam1, CostelloGwilliam2}.
The notation
\[
  \mathfrak A_{K3\times E}
  \;\in\; E_3\text{-}\mathrm{HolFA}(K3\times E)
\]
denotes this chosen finite-stage \(E_3\)-prefactorization datum; the
unconditional functor from CY-\(3\) dg-categories to holomorphic
\(E_3\)-factorization algebras is not asserted
\cite{KontsevichOperadsMotives, TamarkinLittleDisks,
LambrechtsVolicLittleDisks, CostelloGwilliam1, CostelloGwilliam2}.
The retained protected sector supplies a Pfaffian line on the
Dirac--Igusa pro-object.  The Pfaffian--Dirac theorem compares that line
with the automorphic section \(\mathcal D_X=\Delta_5\) only under the named
orientation and wall data.  The five readings are the automorphic
section, the BKM denominator, the retained Pfaffian section, the
conjectural period discriminant, and the singular theta lift.

The unconditional new mathematical content of the manuscript, beyond
the classical Igusa--Borcherds--Gritsenko--Nikulin spine, is the
parity-resolved root-space computation at
\(\delta_{123}=\delta_1+\delta_2+\delta_3\) --- even dimension \(29\)
against odd dimension \(93\), with \(29-93=-64=f(1,1)\)
(Proposition~\ref{prop:bkm-delta123-presentation-count}) --- together
with the scalar-Pfaffian monodromy proposition
(Proposition~\ref{prop:scalar-pfaffian-data-not-section}).  The
Pfaffian--Dirac identity
\(\mathrm{Pf}_{\mathrm{prot}}(\mathfrak D^{DI}_X)=\Delta_5\) is a
recognition criterion relative to the retained data \((D0)\), \((O1)\),
\((O1)^+\), \((O2_{\mathrm{atlas}})\), \((P_{\mathrm{fin}})\), which the
certificate ledgers record as not supplied.

\begin{definition}[Pentadic \(\Delta_5\)]
\label{def:pentadic-delta5}
The five-tuple
\[
  \bigl(\mathcal D_X,\;
  \operatorname{den}(\mathfrak g_{\Delta_5}),\;
  \mathrm{Pf}_{\mathrm{prot}}(\mathfrak D^{DI}_X),\;
  \mathrm{Disc}(\mathcal P\to\mathbb D),\;
  \Theta_{\mathrm{sing}}^{\mathrm{Bor}}(\phi^{\mathrm{Weil}}_{0,1})\bigr)
\]
of automorphic, BKM, retained-Pfaffian, period, and theta-lift readings
is the
\emph{pentadic \(\Delta_5\)}. The reading of \(\Delta_5\) is
fixed by the level: at level \textcircled{1}, \(\Delta_5\) is
\(\mathcal D_X\); at level \textcircled{2}, \(\Delta_5\) is the function whose
value at \(2Z\) is \(64\,\operatorname{den}(\mathfrak g_{\Delta_5})\);
at level \textcircled{3}, \(\Delta_5\) is the protected Pfaffian
\(\mathrm{Pf}_{\mathrm{prot}}(\mathfrak D^{DI}_X)\) under
\((D0)\), \((O1)\), \((O1)^+\), and the retained
\((O2_{\mathrm{atlas}})\) datum; at level
\textcircled{4}, \(\Delta_5^{-2}\) is the mirror discriminant
section, conjecturally; at level \textcircled{5}, \(\Delta_5\) is the
Borcherds theta lift of \(\phi^{\mathrm{Weil}}_{0,1}\).
\end{definition}

\section{The \texorpdfstring{$\textcircled{1}\!\leftrightarrow\!
\textcircled{2}$}{1-2} edge: Borcherds--Gritsenko--Nikulin}
\label{sec:edge-1-2}

The Borcherds product expansion of the weight-zero index-one weak
Jacobi form \(\phi_{0,1}\) is the genus-two cusp form
\[
  \Delta_5(Z)
  \;=\; 64\,q^{1/2}r^{1/2}s^{1/2}
  \prod_{(n,l,m)\in\Gamma_{\mathrm{eff}}}
  \bigl(1 - q^n r^l s^m\bigr)^{f(nm,l)},
\]
where \(Z = \bigl(\begin{smallmatrix} z_1 & z_2 \\ z_2 & z_3 \end{smallmatrix}\bigr)\),
\(q = e^{2\pi i z_1}\), \(r = e^{2\pi i z_2}\), \(s = e^{2\pi i z_3}\),
and \(f(n,l)\) are the Fourier coefficients of \(\phi_{0,1}\) (the
weak Jacobi form whose discriminant decomposition records
\(f(0,0) = 10\), \(f(0,\pm 1) = 1\), \(f(1,0) = 108\),
\(f(1,\pm 1) = -64\), \(f(1,\pm 2) = 10\), \dots)
\cite{Borcherds95, GN, EZ:1}. The leading constant
\([q^{1/2}r^{1/2}s^{1/2}]\Delta_5 = 64 = 2^6\) is the theta-leading
coefficient (Igusa); the chamber
\(0 < |s| \ll |q| \ll |r|^{-1} \ll 1\) fixes the active semigroup
\(\Gamma_{\mathrm{eff}}\) of cusp-completed exponents.

The Gritsenko--Nikulin denominator identity converts the same product
into the Borcherds--Weyl--Kac form on the type-II sublattice
\(\La^{2,1}_{II} \subset \La^{2,1} = \HypPlan\oplus\langle 2\rangle\):
\[
  \operatorname{den}(\mathfrak g_{\Delta_5})
  \;=\; \sum_{w\in W^{(2)}(\La^{2,1}_{II})}
  \mathrm{sgn}(w)\;w\Bigl(e^{\rho}
  \;-\sum_{a\in\La^{2,1}_{II}\cap\,\R_{>0}\Poly_{II}}
  m(a)\,e^{\rho+a}\Bigr)
  \;=\; 64^{-1}\,\Delta_5(2Z),
\]
where \(\Poly_{II}\) is the type-II Weyl chamber (the intersection of
the positive half-spaces of the three simple walls,
Theorem~\ref{thm:bkm-weyl-alternating-sum}) and \(m(a)\in\Z\) are the
imaginary-simple correction multiplicities read from the active support
of \(\phi_{0,1}\) \cite[Theorem~2.1]{GN}, with Weyl vector
\[
  \rho \;=\; \tfrac{1}{2}\bigl(\delta_1 + \delta_2 + \delta_3\bigr)
       \;=\; f_2 - \tfrac{1}{2}\,f_3 + f_{-2}
\]
on the three simple type-II reflections
\(\delta_1 = 2f_2 - f_3\), \(\delta_2 = 2f_{-2} - f_3\),
\(\delta_3 = f_3\) (Cartan matrix
\((\delta_i,\delta_j) = 4\Id_3 - 2\mathbf{J}_3 = \bigl(\begin{smallmatrix}
2 & -2 & -2\\ -2 & 2 & -2\\ -2 & -2 & 2 \end{smallmatrix}\bigr)\),
\(\mathbf{J}_3\) the all-ones matrix); type-II reflection subgroup
\(W^{(2)}(\La^{2,1}_{II})\) generated by these three reflections
\cite{GN, GNPartI, GNII}. The imaginary-simple correction terms
\(-\sum_a m(a)\,e^{\rho+a}\) are the content of the Gritsenko--Nikulin
theorem: the naive alternating sum
\(\sum_w \mathrm{sgn}(w)\,w(e^{\rho})\) alone is the denominator of the
uncorrected hyperbolic Kac--Moody algebra of \(W^{(2)}\) and does not
equal \(64^{-1}\Delta_5(2Z)\) --- the exponent \(\rho + 2f_2\) carries
coefficient \(-9\) in \(64^{-1}\Delta_5(2Z)\) yet lies on no Weyl
translate of \(\rho\), so it must be accounted by imaginary simple
roots. The signed root supermultiplicities recover
the Fourier coefficients of \(\phi_{0,1}\) under the doubling
\(\alpha(n,l,m) = 2nf_2 - lf_3 + 2mf_{-2}\):
\[
  \smult\,\mathfrak g_{\Delta_5,\alpha(n,l,m)}
  \;=\; f(nm,l) \qquad \text{on } (n,l,m)\in\Gamma_{\mathrm{act}}
  =\{(n,l,m)\in\Gamma_{\mathrm{eff}}\mid f(nm,l)\ne 0\}.
\]

The edge \(\textcircled{1}\!\leftrightarrow\!\textcircled{2}\) is the
joint statement: the automorphic section \(\mathcal D_X = \Delta_5\) of
\(\mathcal L^{\otimes 5}\otimes\nu_{\Delta_5}\) on Siegel space
\textsl{equals} the BKM denominator (after the constant \(64^{-1}\) and
the variable doubling \(Z\mapsto 2Z\)).  This is Borcherds 1995
together with the Gritsenko--Nikulin identification of the underlying lattice
with \(\La^{2,1}_{II}\) and the Weyl-group action with
\(W^{(2)}(\La^{2,1}_{II})\) \cite{Borcherds95, GN}.

\section{The \texorpdfstring{$\textcircled{1}\!\leftrightarrow\!
\textcircled{5}$}{1-5} and \texorpdfstring{$\textcircled{2}\!
\leftrightarrow\!\textcircled{5}$}{2-5} edges: the singular theta lift}
\label{sec:edge-1-5}

The Borcherds product expansion of the previous section converts
\(\phi_{0,1}\) into \(\Delta_5\) by an arithmetic identity in
\(\Gamma_{\mathrm{eff}}\); the question of \emph{where} this identity
comes from has a single answer.  The exceptional isomorphism
for \(\PSp_4(\R)\simeq\SO(3,2)_+\) identifies the genus-two Siegel
space with the type-IV symmetric domain attached to the
signature-\((3,2)\) lattice
\(\La^{3,2} = \HypPlan\oplus\HypPlan\oplus\langle 2\rangle\), and it
identifies the arithmetic quotient
\(\Sp_4(\Z)\backslash\mathcal H_2\) with
\(\Gamma_{3,2}\backslash\Hpl^{\IV}_+\), where
\(\Gamma_{3,2}=\wedge^2\Sp_4(\Z)/\{\pm\Id\}\).
Borcherds's regularisation pairs the \(\Mp_2(\Z)\)
Weil-representation form \(\phi^{\mathrm{Weil}}\) with the Siegel
theta kernel of \(\La^{3,2}\), and produces a meromorphic Siegel
modular form whose divisor is
prescribed by the principal part:
\[
  \Theta_{\mathrm{sing}}^{\mathrm{Bor}}(\phi)
  \;=\;
  \int^{\mathrm{reg}}_{\mathrm{SL}_2(\Z)\backslash\Hpl}
  \bigl\langle \phi(\tau),\,
  \Theta_{\La^{3,2}}(\tau, Z)\bigr\rangle\,
  \frac{du\,dv}{v^2},
  \qquad \tau = u + iv,
\]
on the Weil-representation pairing of the vector-valued
weakly-holomorphic input \(\phi\) (the theta decomposition of
\(\phi_{0,1}\) into the discriminant form of \(\La^{3,2}\)) against
the Siegel theta series of the signature-\((3,2)\) lattice
\(\La^{3,2}\) — not the rank-three target lattice \(\La^{2,1}_{II}\),
which receives the BKM root structure under doubling — together with
the antiholomorphic dependence required for absolute convergence after
regularisation; the explicit construction is Chapter~\ref{ch:singular-theta-lift},
Definition~\ref{def:siegel-theta} and
Theorem~\ref{thm:singular-theta-lift} \cite{Borcherds95, B}.  At the
theta-decomposed input \(\phi^{\mathrm{Weil}}_{0,1}\) the regularised
integral computes the composite
\[
  \Theta_{\mathrm{sing}}^{\mathrm{Bor}}\colon
  J^{\mathrm{wk}}_{0,1}
  \;\longrightarrow\;
  M^{!}_{-1/2}(\rho_{\La^{3,2}})
  \;\longrightarrow\;
  M_5\bigl(\Sp_4(\Z),\,\nu_{\Delta_5}\bigr),
  \qquad
  \phi_{0,1} \;\longmapsto\;
  \phi^{\mathrm{Weil}}_{0,1} \;\longmapsto\; \Delta_5,
\]
giving the edge \(\textcircled{1}\!\leftrightarrow\!\textcircled{5}\).
Reading the same lift through the Gritsenko--Nikulin denominator
identity, \(\phi_{0,1}\) is the character of the index-one Heisenberg
sector and the multiplicative lift over \(\La^{2,1}_{II}\) is exactly
\(\operatorname{den}(\mathfrak g_{\Delta_5}) = 64^{-1}\Delta_5(2Z)\),
giving the edge \(\textcircled{2}\!\leftrightarrow\!\textcircled{5}\)
\cite{Borcherds95, GN, GNPartI, GNII}.

\section{The \texorpdfstring{$\textcircled{2}\!\leftrightarrow\!
\textcircled{3}$}{2-3} edge: the Pfaffian--Dirac theorem}
\label{sec:edge-2-3}

The theorem edges proved by Borcherds and Gritsenko--Nikulin identify
the automorphic, BKM-denominator, and singular-theta readings of
\(\Delta_5\); the chiral
\(E_3\)-prefactorization datum of \S\ref{sec:retained-e3-five-readings}
is the retained operator-level source from which the View~\textcircled{3}
datum is obtained.  The structural question has two terms.  The Pfaffian term
asks whether the protected determinant of the \(K3\times E\) operator
datum is the automorphic section \(\Delta_5\).  The primitive term
asks whether the compact Hall primitive sector is the BKM algebra
\(\mathfrak g_{\Delta_5}\).  The first term is the Pfaffian--Dirac
theorem after \((D0)\), \((O1)\), \((O1)^+\), and
\((O2_{\mathrm{atlas}})\), together with \((P_{\mathrm{fin}})\), in
Appendix~\ref{app:obstruction-discharges}.  The second term is the
finite compact-source recognition problem \((S1)\)--\((S10)\).

\begin{theorem}[Pfaffian--Dirac]
\label{thm:pentadic-pfaffian-dirac}
\textsl{Conditional on \((D0)\), \((O1)\), \((O1)^+\),
\((O2_{\mathrm{atlas}})\), and \((P_{\mathrm{fin}})\);
ambient: cosection-twisted reduced d-critical chart on \(K3\times E\)
at finite cusp-truncation height \(R\), with chain-level pro-object,
Mittag--Leffler vanishing of the inverse system in \(R\), and Pfaffian
orientation \(o\) with Weyl character \(\epsilon_o\).} The
protected Pfaffian of the Dirac--Igusa pro-object
\(\mathfrak D^{DI}_X = \varprojlim_R\,\mathfrak D^{DI}_{X,R}\) satisfies
\[
  \mathrm{Pf}_{\mathrm{prot}}(\mathfrak D^{DI}_X)
  \;=\; \Delta_5,
  \qquad
  \epsilon_o
  \;=\; \nu_{\Delta_5}\quad\text{on}\quad
  W^{(2)}(\La^{2,1}_{II}),
\]
\end{theorem}

The primitive Lie-superalgebra assertion is a separate recognition
statement.  If the finite compact-source datum \((S1)\)--\((S10)\) is
supplied on a cofinal system of relation-closed degree windows, then
\[
  \mathcal P_X
  :=
  H^0\,\Prim_{\mathrm{prot}}\bigl(\overline\Pi^{\Theta}_{X,*}
  \mathcal F_X\bigr) /
  \overline{\operatorname{Rad}\langle\,,\rangle_X}
  \;\cong\; \mathfrak g_{\Delta_5}.
\]

The four obstructions and the finite geometric Pfaffian datum name the
exact conditioning. \((D0)\) is the
\(D0\)-degeneration limit of the cosection-reduced d-critical
orientation theory: the limiting orientation in the \(D0\)-degenerate
fibre is constructed and matches the orientation of the generic fibre
under the cosection-twist
\cite{BBDJS2015, JoyceUpmeier, JoyceTanakaUpmeier}. \((O1)\) is the
strong reduced orientation on the \(K3\times E\)-quotient self-Ext
frame bundle: the square root of the determinant line of the cotangent
complex extends across the cusp boundary as part of the retained
quotient-orientation datum
\cite{CaoGrossJoyce, JoyceUpmeier, JoyceTanakaUpmeier}. \((O1)^+\) is the Weyl-equivariant
transport of the orientation along reflections in the type-II
reflection subgroup \(W^{(2)}(\La^{2,1}_{II})\): the orientation cocycle
\(c_o\in Z^2(W^{(2)}(\La^{2,1}_{II});\mathbb F_2)\) trivializes in
\(H^2\), and the resulting character
\(\epsilon_o:W^{(2)}(\La^{2,1}_{II})\to\{\pm1\}\) is well-defined on
generators. \((O2)\) is the local Pfaffian wall normal form on type-II
walls: at each simple type-II wall \(s_{\delta_i}\) the rank-one normal
Pfaffian crossing of Definition~\ref{def:rank-one-type-ii-crossing}
is part of the retained wall-atlas datum and gives the local sign
\((-1)^{N^{\mathrm{Pf}}_{\delta_i}} = -1\), and the Hall-side sign
sum on the half-Hilbert orbit of the self-Ext stack of size
\(2^{12}=4096\) is compared with the squared theta-leading
\(([q^{1/2}r^{1/2}s^{1/2}]\Delta_5)^2\) only after scalarization.  The
wall transport gives the character value
\(\epsilon_o(s_{\delta_i}) = \nu_{\Delta_5}(s_{\delta_i})\) on each
generator (Appendix~E).
The finite geometric Pfaffian datum gives the skew perfect complexes,
Pfaffian lines, finite Pfaffian sections, type-II wall divisor orders,
automorphic line comparisons, support and parity rows, leading
coefficient, and Pfaffian-line Mittag--Leffler compatibility for the
cofinal HN product.

\begin{remark}[The conditioning is named precisely]
\label{rem:conditioning-named}
The four obstruction classes \((D0)\), \((O1)\), \((O1)^+\), \((O2)\)
and the finite geometric Pfaffian datum are the exact Pfaffian and
orientation obligations; Appendix~\ref{app:obstruction-discharges}
reduces \((D0)\) to the retained D0-HN datum, names the retained
quotient-orientation and Weyl transport data, and names
\((O2_{\mathrm{atlas}})\) and \((P_{\mathrm{fin}})\). At the operator level
\textcircled{3}, \(\mathcal D_X=\Delta_5\) is compared with the protected
Pfaffian \(\mathrm{Pf}_{\mathrm{prot}}(\mathfrak D^{DI}_X)\).  After
the level-\(n\) protected trace, the scalar partition function is
\(\Delta_5^{-2}\). The primitive algebra
\(\mathcal P_X\simeq\mathfrak g_{\Delta_5}\) is stronger; it requires
the finite compact-source recognition datum \((S1)\)--\((S10)\). The
operator-level datum carries the orientation character
\(\epsilon_o\) that the scalar shadow has squared away, by view
\textcircled{3}, and recovers it as \(\nu_{\Delta_5}\) on
\(W^{(2)}(\La^{2,1}_{II})\) by \((O1)^+\) and \((O2)\).
\end{remark}

\section{View \texorpdfstring{$\textcircled{4}$}{4}: mirror
discriminant}
\label{sec:view-4-mirror}

\textsl{Conjectural} (ambient: variation of Hodge structure on the
rank-\(3\) Mukai-invariant \((K3\times E)\)-period component).
On \(\mathcal P_{K3\times E}^{(3)}\simeq \mathcal H_2/\Sp_4(\Z)\), the
mirror discriminant section is conjecturally \(\Delta_5^{-2}\).  Its
polar divisor is \(2\mathfrak D_{\mathrm{mir}}\), where the reduced
support of \(\mathfrak D_{\mathrm{mir}}\) is the Humbert--Heegner
divisor \(\mathfrak H_{\mathrm{II}}\) fixed by the Borcherds product of
\(\Delta_5\).  This mirror-discriminant identification is conjectural
and is not a hypothesis of any
theorem; see Conjecture~\ref{conj:mirror-discriminant}.

\section{The level-\texorpdfstring{$n$}{n} scalar shadow}
\label{sec:level-n-scalar-shadow}

The scalar partition function on the closed cobordism
\(K3\times E\mapsto\mathrm{pt}\) is the Borcherds determinant of the
protected K3 index, computed by Oberdieck--Pixton 2018 and
Oberdieck--Pandharipande 2018 \cite{OB, OPand}: the OP partition
function \(Z^X_{\mathrm{OP}} = -4096\,\Delta_5^{-2}\) is the virtual
count of Hilbert schemes of \(K3\times E\) in the
\((-p_{\mathrm{DT}})^n\) DT-chamber, and the BPS normalization
\(Z^{K3\times E}_{\mathrm{BPS}} = \Delta_5^{-2}\) rescales by
\(-4096^{-1}\).
Two normalizations of the same weight-ten cusp form on \(\Sp_4(\Z)\)
are in use, and naming them is necessary because both appear in the
formulas:
\begin{itemize}
\item \emph{Theta-normalized.}
\(\Delta_{10}(Z) := \Delta_5(Z)^2 \in M_{10}(\Sp_4(\Z))\) (the
multiplier \(\nu_{\Delta_5}\) squares to trivial; the leading
coefficient \([q\,r\,s]\Delta_{10} = 64^2 = 4096\)).
\item \emph{Oberdieck--Pixton-monic.}
\(\chi_{10}^{\mathrm{OP}}(Z) := D_5(Z)^2 = 64^{-2}\Delta_5(Z)^2 =
4096^{-1}\Delta_{10}(Z)\) (so \([q\,r\,s]\chi_{10}^{\mathrm{OP}} = 1\)).
\end{itemize}
Oberdieck--Pixton and Oberdieck--Pandharipande prove
\cite{OB, OPand}
\[
  Z^X_{\mathrm{OP}}
  \;=\; -\bigl(\chi_{10}^{\mathrm{OP}}\bigr)^{-1}
  \;=\; -4096\,\Delta_5^{-2},
\]
and the level-\(n\) protected scalar shadow is
\[
  Z^{K3\times E}_{\mathrm{BPS}}
  \;=\; \bigl(\Delta_{10}\bigr)^{-1}
  \;=\; \Delta_5^{-2}
  \;=\; -\tfrac{1}{4096}\,Z^X_{\mathrm{OP}}.
\]
The factor \(4096 = 64^2\) is the squared theta-leading constant of
\(\Delta_5\); the OP minus sign is a scalar branch convention absorbed
into the OP \((-p_{\mathrm{DT}})^n\) chamber and is independent of the
orientation character \(\epsilon_o\) of view~\textcircled{3}.

\(Z^{K3\times E}_{\mathrm{BPS}}\) is the level-\(n\) protected trace of
the retained \(K3\times E\) data: a charge-graded virtual partition
function.  A three-dimensional gravity-line interpretation of the
scalar shadow with internal compact factor \(K3\times E\) requires the
Hall--Borcherds residual of Vol II.  The universal-stage
chain names the required objects: primitive factorization dg-category
\(\mathcal C\) → boundary chart \(A_b = \End_{\mathcal C}(b)\) → bar
twisting \(\mathrm{Bar}(A_b)\) → derived chiral centre
\(Z^{\mathrm{der}}_{\mathrm{ch}}(A_b)\) → ambient operator algebra
\(\mathcal A\) at the boundary vacuum \(b\) of the open
factorization dg-category on
\((K3\times E,\,D,\,\tau)\); this chain is not constructed here.
The local witness is the formal-Darboux stalk of the
mixed-holomorphic-topological strings theory
\cite{KontsevichOperadsMotives, CostelloLiBCOV,
CostelloLiOpenClosedBCOV}.  The class
\(\mathrm{ob}^{\mathrm{hol}}_{\mathrm{dR}}(K3\times E)\) is the
holomorphic de Rham obstruction to gluing this local BV factorization
model over \(K3\times E\); its vanishing is a necessary global
condition, not the gravity-line operator algebra.

\begin{remark}[Primitive order]
\label{rem:beilinson-cut-delta5}
The five views are ordered: primitive objects (the retained chiral
\(E_3\)-prefactorization datum on \(K3\times E\), its four obstructions
\((D0),(O1),(O1)^+,(O2_{\mathrm{atlas}})\), and its finite Pfaffian
section datum; the Cartan datum of \(\mathfrak g_{\Delta_5}\))
first; cross-level identities ((②↔③) Pfaffian--Dirac, (\(\epsilon_o
=\nu_{\Delta_5}\)) orientation match, \((\mathcal P_X\cong\mathfrak
g_{\Delta_5})\) primitive recognition) second; scalar modular forms
last (the section \(\mathcal D_X\), the BKM denominator
\(64^{-1}\Delta_5(2Z)\), the partition function \(\Delta_5^{-2}\)).
A statement is not primitive if it is true only after passing to a
trace, choosing a chamber, completing a category, or installing
descent data.  On each first appearance the scalar reading is named:
the theta-normalized form \(\Delta_5\) with leading
\([q^{1/2}r^{1/2}s^{1/2}]\Delta_5 = 64\); the BKM denominator
\(64^{-1}\Delta_5(2Z)\) with the doubling \(2Z\) on the type-II
sublattice; the OP-monic form \(D_5 = 64^{-1}\Delta_5\) with
\(\chi_{10}^{\mathrm{OP}} = D_5^2 = 4096^{-1}\Delta_5^2\) the
Oberdieck--Pixton normalization of the squared determinant section.
\end{remark}

\section{Placement in the Beilinson Chain}
\label{sec:beilinson-chain-placement}

The pentadic \(\Delta_5\) is one identification at one specific level
of the Beilinson chain. Under the
chiral Koszul source datum of
Definition~\ref{def:chiral-koszul-source-datum}, the finite Koszul
comparison datum of
Definition~\ref{def:finite-koszul-comparison-datum}, and the strict
Koszul Mittag--Leffler exactness of
Proposition~\ref{prop:chiral-koszul-source-datum-consequence}, the
source coalgebra \(C_X\) maps by cobar comparison to the target chiral
shadow \(\mathsf{Den}_{\Delta_5,E}\), compatibly with Weyl transport,
Pfaffian orientation, Hall pairing, radical quotient, and PBW
filtration.  Vol I's Theorem H supplies the
chain-level identification
\(Z^{\mathrm{der}}_{\mathrm{ch}}(A_b) =
\mathrm{ChirHoch}^\bullet(A_b,A_b)\) only on the retained Koszul
locus, and Vol II's chiral-Hochschild Beilinson stratification
\textnormal{(i)+(ii)} supplies the CY-orientation trace pairing of
degree \(-\dim_{\mathbb C}(K3\times E) = -3\).  Relative to the
retained trace datum, these inputs give the level-\(\mathsf Z\) trace
identity
\(\mathrm{Tr}^{\mathrm{bulk}}_n((-1)^F\cdot\mathrm{id})_{Z^{\mathrm{der}}_{\mathrm{ch}}(A_b)}
	= \Delta_5^{-2}\) of Appendix~\ref{app:obstruction-discharges},
Theorem~\ref{thm:G-vol2-discharge}, relative to \((Z_{\mathrm{HB}})\);
the OP chamber scalar is
\(-4096\) times this trace. The level-\(\mathsf A\)
gravity-line operator algebra acting on the boundary is the open
gravity-line Hall--Borcherds residual of Vol II.  The Vol III
\(\kappa\)-stratification places
\(\mathfrak g_{\Delta_5}\) at the \((d=3,\,\text{rank }3)\) chamber
through the identity \(\kappa_{\mathrm{BKM}}(\Phi_N) = c_N(0)/2\),
where \(c_N(0)\) is the constant Fourier coefficient of the input weak
Jacobi or vector-valued modular form. Applied at the paramodular row
\((t,N;k) = (1,1;5)\) of the Gritsenko--Cl\'ery catalogue
\cite{GC, GN}, with Jacobi seed \(\phi_{0,1}\) and Borcherds product
\(\Delta_5\), the identity reads
\(\kappa_{\mathrm{BKM}}(\Delta_5) = c_1(0)/2 = 5\).  The companion
identity at input \(\Phi_{12}\), the Fake-Monster Borcherds product
on the Leech-extension lattice \(\mathrm{II}_{25,1}\)
\cite{BorcherdsGKM88, Borcherds95}, reads
\(\kappa_{\mathrm{BKM}}(\Phi_{12}) = c_{\Phi_{12}}(0)/2 = 12\); the
two values are different rows of one universal \(c_N(0)/2\) rule,
indexed by different input forms and lattices, not in contradiction.
The mixed holomorphic-topological strings companion identifies the
formal-Darboux stalk
\(A^{\mathrm{cl}}_{\partial,N} = C^\bullet_{\mathrm{CE}}(\mathfrak{gl}_N,
\,\mathrm{Kosz}([\phi_1,\phi_2]))\) of the chiral Hochschild cochain
at the Dirac brane vacuum.  The obstruction to globalizing this local
factorization model over \(K3\times E\) is the holomorphic de Rham
class \(\mathrm{ob}^{\mathrm{hol}}_{\mathrm{dR}}(K3\times E)\).  A
three-dimensional gravitational path-integral interpretation of the
level-\(n\) scalar shadow
\(Z^{K3\times E}_{\mathrm{BPS}} = \Delta_5^{-2}\) requires this
obstruction to vanish and separately requires the Vol II gravity-line
Hall--Borcherds operator algebra.
